\section{Limitations and Future Work}

A key limitation of our current approach is the difficulty of training a single shared residual policy across multiple demonstrations simultaneously. While training on individual demonstrations yields reliable manipulation policies, extending to multiple demonstrations introduces several challenges. First, distributing a fixed number of parallel simulation environments across many demonstrations reduces the per-demonstration sample count, which can destabilize policy optimization. Second, and more fundamentally, demonstrations of differing object geometries, contact configurations, or motion trajectories may induce conflicting policy gradients—corrections beneficial for one demonstration may be detrimental to another—causing the shared policy to converge to a compromise that satisfies one demonstration rather than another. Additionally, there are limitations to learning, as there is noise from the AVP and OptiTrack systems; specifically, the Apple Vision struggles when there are occlusions of the hand due to object interactions. A further limitation stems from inaccuracies in the object meshes and small errors in their estimated poses. Such mismatches caused failures in simulation — for instance, collisions between meshes — despite the corresponding real-world configuration being properly aligned.

\section{Conclusion}

We introduce a bi-manual teleoperation system that takes a concrete step toward generalizable dexterous manipulation. Built on our proposed multi-demonstration residual policy with cross-system spatial calibration, our approach learns contact-aware manipulation without task-specific reward engineering and operates on unseen operator trajectories without retraining. Across two long-horizon bi-manual tasks, our method is the only approach to complete the full manipulation sequence end to end, where both the pre-trained imitator and optimization-based retargeting uniformly fail at contact-critical phases. Furthermore, the same hardware pipeline that collects training demonstrations drives live teleoperation without modification, establishing a scalable foundation for dexterous bi-manual data collection and a viable path toward robust human-level robotic manipulation.

\section*{Acknowledgments}

The author would like to thank Cheng Pan and Prof.\ Josie Hughes for their guidance and support throughout this project. Parts of the code and report were developed with the assistance of Claude (Anthropic), used for code generation, debugging, and drafting of written content. All outputs were reviewed, verified, and edited by the author.